\begin{proof}[Proof of \cref{obj:lem:bounded-lower-pair}]
\begin{enumerate}
\item We first obtain a Hölder radius for the scaled bump that is uniform in its scale and center. Set
\[
m:=\max\{\lceil\beta\rceil-1,0\},
\qquad
\sigma_\beta:=\beta-m\in(0,1].
\]
For each derivative order \(\kappa\in\mathbb N\), write
\[
\mathscr J_\kappa(z):=\mathrm d^\kappa\psi(z),
\qquad z\in\mathbb R^d,
\qquad
\mathscr J_0(z)=\psi(z),
\]
where \(\mathrm d^\kappa\) denotes the \(\kappa\)-th Fréchet derivative, with the operator norm induced by the supremum norm.

The coordinate factors defining \(\psi\) are smooth across their endpoints: their derivatives from inside \((-1,1)\) vanish at the endpoints because the exponential factor decays faster than every inverse power of the distance to the endpoint. Consequently,
\[
\psi\in C^\infty(\mathbb R^d),
\qquad
\operatorname{supp}\psi\subseteq[-1,1]^d.
\]
Every \(\mathscr J_\kappa\) is continuous and compactly supported. Thus the constants
\[
M_{\psi,\kappa}:=
\sup_{z\in\mathbb R^d}\|\mathscr J_\kappa(z)\|,
\qquad \kappa\in\mathbb N,
\]
are finite and nonnegative. Define
\[
M_{\psi,*}:=\max\{M_{\psi,m},M_{\psi,m+1}\},
\qquad
M_{\psi,\mathrm{mod}}:=2M_{\psi,*}.
\]
The derivative bound of order \(m+1\), applied along a line segment, gives
\[
\|\mathscr J_m(z)-\mathscr J_m(z')\|
\le M_{\psi,*}\|z-z'\|_\infty.
\]
If \(\|z-z'\|_\infty\le1\), then
\[
\|\mathscr J_m(z)-\mathscr J_m(z')\|
\le M_{\psi,*}\|z-z'\|_\infty^{\sigma_\beta}
\le M_{\psi,\mathrm{mod}}\|z-z'\|_\infty^{\sigma_\beta}.
\]
If \(\|z-z'\|_\infty>1\), then
\[
\|\mathscr J_m(z)-\mathscr J_m(z')\|
\le 2M_{\psi,*}
\le M_{\psi,\mathrm{mod}}\|z-z'\|_\infty^{\sigma_\beta}.
\]
These two cases include \(\sigma_\beta=1\).

For \(0<h\le1\), define the scaled bump by
\[
\psi_h(x):=h^\beta\psi((x-x_0)/h),
\qquad x\in\mathbb R^d.
\]
Translation and dilation give, for \(0\le\kappa\le m\),
\[
\mathrm d^\kappa\psi_h(x)
=h^\beta h^{-\kappa}\mathscr J_\kappa((x-x_0)/h),
\]
and hence
\[
\|\mathrm d^\kappa\psi_h(x)\|
\le h^{\beta-\kappa}M_{\psi,\kappa}
\le M_{\psi,\kappa}.
\]
For the top derivative, the preceding modulus bound yields
\[
\begin{aligned}
\|\mathrm d^m\psi_h(x)-\mathrm d^m\psi_h(x')\|
&\le h^{\beta-m}M_{\psi,\mathrm{mod}}
       \left(\frac{\|x-x'\|_\infty}{h}\right)^{\sigma_\beta}\\
&=M_{\psi,\mathrm{mod}}\|x-x'\|_\infty^{\sigma_\beta},
\end{aligned}
\]
since \(\beta-m=\sigma_\beta\). Set
\[
A_\psi:=
\max\left\{
1,\,
\max\left\{
\sum_{\kappa=0}^{m}M_{\psi,\kappa},
M_{\psi,\mathrm{mod}}
\right\}
\right\}>0.
\]
Coordinate derivatives are evaluations of these multilinear derivatives on coordinate unit vectors, so their magnitudes and differences are bounded by the corresponding operator norms. Restricting the ambient smooth functions to the cube supplies their continuous boundary traces. Therefore, with the convention of \cref{obj:def:holder},
\[
\psi_h|_{\mathcal X}\in\mathcal H^\beta(A_\psi)
\qquad
\text{for every }0<h\le1,
\quad \mathcal X=[0,1]^d.
\]
% lean: Causalean.Mathlib.Analysis.Calculus.CubeExtension.exists_uniform_scaledProductBump_holderBallOn

\item Using the baseline probability from \cref{obj:def:bounded-lower-pair}, choose
\[
\theta_0:=\min\left\{\frac14,\frac{L}{4B}\right\},
\qquad
a:=\min\left\{\frac{\theta_0}{2},\frac{L}{4BA_\psi}\right\}.
\]
Since \(B,L,A_\psi>0\),
\[
0<\theta_0\le\frac14,
\qquad
0<a\le\theta_0.
\]
Define the effective dimension and entropy coefficient by
\[
D_\gamma:=\frac{d\gamma}{\gamma-1},
\qquad
\tau_{\mathrm{KL}}:=2\beta+D_\gamma,
\qquad
K_{\mathrm{KL}}:=\frac{2}{\theta_0}\,2^{D_\gamma}.
\]
Here \(D_\gamma>0\), \(K_{\mathrm{KL}}>0\), and
\(\tau_{\mathrm{KL}}\ge1\). Choose
\[
c:=\frac{1}{8(K_{\mathrm{KL}}a^2+1)},
\qquad
c_0:=Bac^\beta.
\]
Then \(0<c\le1\) and \(c_0>0\). Because \(c\le1\) and
\(\tau_{\mathrm{KL}}\ge1\),
\[
c^{\tau_{\mathrm{KL}}}\le c,
\]
so
\[
K_{\mathrm{KL}}a^2c^{\tau_{\mathrm{KL}}}
\le K_{\mathrm{KL}}a^2c
=\frac{K_{\mathrm{KL}}a^2}{8(K_{\mathrm{KL}}a^2+1)}
\le\frac18.
\]
All these constants are fixed independently of \(n\).
% lean: CausalSmith.Stat.WeakOverlap.lowerPair_choose_kl_scale

\item Fix an integer \(n\ge1\). The oracle mesh and the selected bump scale are
\[
h_{\gamma,n}:=n^{-1/\tau_{\mathrm{KL}}},
\qquad
h:=ch_{\gamma,n}.
\]
Since \(n\ge1\), \(\tau_{\mathrm{KL}}>0\), and \(0<c\le1\),
\[
0<h_{\gamma,n}\le1,
\qquad
0<h\le1.
\]
Use \(R,F_R,e_0,\psi,p_{j,h}\) and \(P_{j,h}\) as defined in
\cref{obj:def:bounded-lower-pair}. Every factor defining \(F_R\) lies in
\([0,1]\), and the exponent \(1/(\gamma-1)\) is positive. Thus
\[
0\le F_R(R(x))\le1,
\qquad
0\le e_0(x)\le1,
\qquad x\in\mathbb R^d.
\]
For each coordinate factor of \(\psi\), the exponent is nonpositive when
the coordinate has absolute value less than one, and the factor is zero
otherwise. Hence
\[
0\le\psi(z)\le1,
\qquad z\in\mathbb R^d.
\]
Together with \(h^\beta\le1\), this gives
\[
0\le ah^\beta\psi((x-x_0)/h)\le a
\]
and therefore
\[
0<\theta_0\le p_{j,h}(x)
\le\theta_0+a
\le2\theta_0
\le\frac12,
\qquad j\in\{0,1\},\quad x\in\mathbb R^d.
\]
In particular, \(a,c\) lie in the admissible domain of
\cref{obj:def:bounded-lower-pair}, simultaneously for every \(n\ge1\).

Write the constructed completion laws as
\[
\widetilde P_{j,h}:=
\mathcal L(X,A,Y(0),Y(1)),
\qquad j\in\{0,1\},
\]
where the joint law is the product-kernel construction in
\cref{obj:def:bounded-lower-pair}. The cube has Lebesgue volume
\[
\operatorname{vol}(\mathcal X)=\prod_{i=1}^{d}(1-0)=1.
\]
The admissible weights sum to one in each two-point factor. The kernels
are measurable, so integrating them over the cube gives probability
laws \(\widetilde P_{j,h}\), whose observed images are \(P_{j,h}\).
% lean: CausalSmith.Stat.WeakOverlap.boundedLowerPairAdmissible_of_small

\item Define the displayed treated-response curves by
\[
\mu_{j,h}(x):=Bp_{j,h}(x),
\qquad x\in\mathbb R^d,\quad j\in\{0,1\}.
\]
The choices of \(\theta_0\) and \(a\) give
\[
|B\theta_0|=B\theta_0\le\frac L4,
\qquad
|Ba|A_\psi=BaA_\psi\le\frac L4,
\]
and consequently
\[
|B\theta_0|+|Ba|A_\psi\le L.
\]
For the baseline curve, all positive-order derivatives vanish and
\[
\|\mu_{0,h}\|_\infty=B\theta_0\le L.
\]
For the perturbed curve,
\[
\mu_{1,h}(x)=B\theta_0+Ba\psi_h(x).
\]
For every multi-index \(\alpha\) with \(|\alpha|\le m\),
\[
D^\alpha\mu_{1,h}(x)
=
\begin{cases}
B\theta_0+Ba\psi_h(x),&|\alpha|=0,\\
BaD^\alpha\psi_h(x),&|\alpha|>0.
\end{cases}
\]
The bounds from the first step imply
\[
\|D^\alpha\mu_{1,h}\|_\infty
\le |B\theta_0|+|Ba|A_\psi
\le L.
\]
For \(|\alpha|=m\), the constant term cancels in the difference, giving
\[
|D^\alpha\mu_{1,h}(x)-D^\alpha\mu_{1,h}(x')|
\le |Ba|A_\psi\|x-x'\|_\infty^{\sigma_\beta}
\le L\|x-x'\|_\infty^{\sigma_\beta}.
\]
Their ambient smoothness supplies the required boundary traces.
It follows from \cref{obj:def:holder} that
\[
\mu_{0,h},\mu_{1,h}\in\mathcal H^\beta(L).
\]

The treated potential-outcome marginal of the conditional product kernel
is \(B\operatorname{Bernoulli}(p_{j,h}(x))\). Thus, under the constructed
completion,
\[
\mathbb E_{\widetilde P_{j,h}}[Y(1)\mid X=x]
=(1-p_{j,h}(x))\,0+p_{j,h}(x)B
=\mu_{j,h}(x)
\]
for Lebesgue-almost every \(x\in\mathcal X\). This establishes the
response smoothness requirement in \cref{obj:ass:holder-response}.
% lean: CausalSmith.Stat.WeakOverlap.pairSuccess_true_holderBall
% lean: CausalSmith.Stat.WeakOverlap.completionBind_holderResponse

\item We next establish the propensity tail bound. For every \(r\ge0\),
the intersection of the radial sublevel set with the cube is
\[
\{x:R(x)\le r\}\cap\mathcal X
=
\prod_{i=1}^{d}
\left[
\max\{0,x_{0i}-r\},
\min\{1,x_{0i}+r\}
\right].
\]
Since \(x_0\in\mathcal X\), these intervals are nonempty. Their product
volume gives
\[
\mathbb P\{R(X)\le r\}=F_R(r)
\]
under uniform \(X\). The displayed formula for \(F_R\) is continuous,
and the nested sublevel sets show that it is nondecreasing on
\([0,\infty)\). Also,
\[
F_R(0)=0,\qquad F_R(1)=1.
\]
The first identity uses \(d>0\); the second follows because every
coordinate interval at radius one is \([0,1]\).

For \(u_{\mathrm{tail}}\in[0,1)\), define
\[
\mathscr R(u_{\mathrm{tail}})
:=
\{r\in[0,\infty):F_R(r)\le u_{\mathrm{tail}}\},
\qquad
r_{\sup}(u_{\mathrm{tail}})
:=\sup\mathscr R(u_{\mathrm{tail}}).
\]
The set contains zero, is closed by continuity, and is bounded above
by one because \(F_R(r)\ge F_R(1)=1\) for \(r\ge1\). Its supremum
therefore belongs to the set. Consequently,
\[
\begin{aligned}
\mathbb P\{F_R(R(X))\le u_{\mathrm{tail}}\}
&\le \mathbb P\{R(X)\le r_{\sup}(u_{\mathrm{tail}})\}\\
&=F_R(r_{\sup}(u_{\mathrm{tail}}))
\le u_{\mathrm{tail}}.
\end{aligned}
\]
For \(u_{\mathrm{tail}}=1\), the same upper bound follows from total
probability one.

For \(t\in[0,1]\), raising \(e_0(x)\le t\) to the positive power
\(\gamma-1\) gives
\[
\{x:e_0(x)\le t\}
\subseteq
\{x:F_R(R(x))\le t^{\gamma-1}\}.
\]
Since \(t^{\gamma-1}\in[0,1]\), the preceding bound yields
\[
\mathbb P\{e_0(X)\le t\}
\le t^{\gamma-1}
\le Ct^{\gamma-1}.
\]
At \(t=0\), this shows
\[
\mathbb P\{e_0(X)=0\}=0.
\]
Together with \(e_0\ge0\), it follows that \(e_0(X)>0\) almost surely.
% lean: CausalSmith.Stat.WeakOverlap.radialCDF_self_sublevel_mass
% lean: CausalSmith.Stat.WeakOverlap.radialPropensity_uniform_tail
% lean: CausalSmith.Stat.WeakOverlap.radialPropensity_uniform_pos_ae

\item We verify the remaining model requirements for the constructed
completions. Their covariate marginal is uniform on the cube, so
\[
f(x)=1\ge c_f
\qquad\text{for Lebesgue-almost every }x\in\mathcal X.
\]
The conditional product kernel has the factorization
\[
\begin{aligned}
&\mathcal L_{\widetilde P_{j,h}}
   \bigl((Y(0),Y(1)),A\mid X=x\bigr)\\
&\qquad=
\left(
B\operatorname{Bernoulli}(\theta_0)
\otimes
B\operatorname{Bernoulli}(p_{j,h}(x))
\right)
\otimes\operatorname{Bernoulli}(e_0(x)).
\end{aligned}
\]
Thus \((Y(0),Y(1))\perp A\mid X\), and
\[
\mathbb P_{\widetilde P_{j,h}}\{A=1\mid X=x\}=e_0(x)
\]
almost everywhere. The observation map enforces
\[
Y=AY(1)+(1-A)Y(0)
\qquad\widetilde P_{j,h}\text{-almost surely}.
\]

The endpoint-supported potential-outcome kernels and this observation
map imply \(0\le Y\le B\) almost surely. To apply this support bound to
the treated conditional distribution, define
\[
\nu_{j,x}:=\mathcal L_{P_{j,h}}(Y\mid X=x,A=1).
\]
Disintegrating the support statement first over \(X\) and then over
\(A\), and using \(e_0(X)>0\) almost surely, gives
\[
\nu_{j,x}([0,B])=1
\qquad\text{for Lebesgue-almost every }x\in\mathcal X.
\]
For such \(x\), put
\[
\overline y_{j,x}:=\int y\,\nu_{j,x}(dy).
\]
Integrating the support bounds gives
\[
0\le\overline y_{j,x}\le B.
\]
The bounded-variable sub-Gaussian moment bound gives, for every
\(\lambda\in\mathbb R\),
\[
\int
\exp\{\lambda(y-\overline y_{j,x})\}\,\nu_{j,x}(dy)
\le
\exp\left\{\frac{B^2\lambda^2}{8}\right\}
\le
\exp\left\{\frac{B^2\lambda^2}{2}\right\}.
\]
The exponential is integrable because \(y\in[0,B]\) almost surely.
These bounds establish
\cref{obj:ass:bounded-potential-outcomes}. The density, consistency,
exchangeability, propensity-tail, and smoothness properties proved
above establish the other requirements of \cref{obj:def:model}.
Therefore, with these exact completion witnesses and response curves,
\[
P_{0,h},P_{1,h}\in\mathcal P_\gamma.
\]
% lean: CausalSmith.Stat.WeakOverlap.completionBind_uniformCube_globalTailModel
% lean: CausalSmith.Stat.WeakOverlap.completionBind_boundedOutcome_of_endpoint_support

\item For completeness, the two potential-outcome factors at every
covariate point assign total mass one to \(\{0,B\}\). Integrating their
product over the cube therefore gives
\[
\widetilde P_{j,h}
\{Y(0)\in\{0,B\},\ Y(1)\in\{0,B\}\}
=
\int_{\mathcal X}1\,dx
=1,
\qquad j\in\{0,1\}.
\]
This is the required endpoint support under each constructed completion.
% lean: CausalSmith.Stat.WeakOverlap.lowerPairCompletion_supported_on_endpoints

\item We now bound the entropy of the conditional observation laws.
For \(x\in\mathbb R^d\), define
\[
p_{\mathrm{Ber},x}:=p_{0,h}(x),
\qquad
q_{\mathrm{Ber},x}:=p_{1,h}(x).
\]
Both parameters lie in \((0,1)\), and
\[
\theta_0\le q_{\mathrm{Ber},x}\le\frac12,
\qquad
q_{\mathrm{Ber},x}(1-q_{\mathrm{Ber},x})
\ge\frac{\theta_0}{2}.
\]
Applying \(\log v\le v-1\), for \(v>0\), to each likelihood-ratio term gives
\[
\begin{aligned}
&\operatorname{KL}\!\left(
\operatorname{Bernoulli}(p_{\mathrm{Ber},x}),
\operatorname{Bernoulli}(q_{\mathrm{Ber},x})
\right)\\
&\quad=
p_{\mathrm{Ber},x}\log\frac{p_{\mathrm{Ber},x}}{q_{\mathrm{Ber},x}}
+(1-p_{\mathrm{Ber},x})
 \log\frac{1-p_{\mathrm{Ber},x}}{1-q_{\mathrm{Ber},x}}\\
&\quad\le
p_{\mathrm{Ber},x}
 \left(\frac{p_{\mathrm{Ber},x}}{q_{\mathrm{Ber},x}}-1\right)
+(1-p_{\mathrm{Ber},x})
 \left(\frac{1-p_{\mathrm{Ber},x}}{1-q_{\mathrm{Ber},x}}-1\right)\\
&\quad=
\frac{(p_{\mathrm{Ber},x}-q_{\mathrm{Ber},x})^2}
     {q_{\mathrm{Ber},x}(1-q_{\mathrm{Ber},x})}\\
&\quad\le
\frac{2}{\theta_0}(p_{\mathrm{Ber},x}-q_{\mathrm{Ber},x})^2.
\end{aligned}
\]
Also,
\[
0\le p_{1,h}(x)-p_{0,h}(x)
=ah^\beta\psi((x-x_0)/h)
\le ah^\beta.
\]

Define the laws of the treatment and binary selected outcome by
\[
\mathsf V_{j,x}:=
(1-e_0(x))\,\delta_0\otimes\operatorname{Bernoulli}(\theta_0)
+
e_0(x)\,\delta_1\otimes\operatorname{Bernoulli}(p_{j,h}(x)),
\]
on \(\{0,1\}^2\). Their scaled and recorded versions are
\[
\mathsf W_{j,x}
:=
\bigl((b_A,b_Y)\mapsto(b_A,Bb_Y)\bigr)_\#\mathsf V_{j,x},
\qquad
\mathsf Q_{j,x}
:=
\bigl((b_A,y)\mapsto(x,b_A,y)\bigr)_\#\mathsf W_{j,x}.
\]
Here the first map has domain \(\{0,1\}^2\), the second has domain
\(\{0,1\}\times\mathbb R\), and \(\#\) denotes image measure.
The product construction shows that \(\mathsf Q_{j,x}\) is the
conditional observation law at covariate \(x\).

The control parameters agree, so the control contribution to entropy
is zero. The treated contribution satisfies the preceding quadratic
bound. Conditioning on the binary treatment therefore gives
\[
\begin{aligned}
\operatorname{KL}(\mathsf V_{0,x},\mathsf V_{1,x})
&=(1-e_0(x))
  \operatorname{KL}\!\left(
  \operatorname{Bernoulli}(\theta_0),
  \operatorname{Bernoulli}(\theta_0)\right)\\
&\quad+
e_0(x)\operatorname{KL}\!\left(
\operatorname{Bernoulli}(p_{0,h}(x)),
\operatorname{Bernoulli}(p_{1,h}(x))\right)\\
&\le \frac{2}{\theta_0}(ah^\beta)^2e_0(x).
\end{aligned}
\]
Applying entropy contraction first to recording the covariate and then
to scaling the binary outcome yields
\[
\operatorname{KL}(\mathsf Q_{0,x},\mathsf Q_{1,x})
\le
\operatorname{KL}(\mathsf W_{0,x},\mathsf W_{1,x})
\le
\operatorname{KL}(\mathsf V_{0,x},\mathsf V_{1,x})
\le
\frac{2}{\theta_0}(ah^\beta)^2e_0(x).
\]
In particular, the conditional entropies are finite.
% lean: CausalSmith.Stat.WeakOverlap.lowerPair_bernoulli_kl_le_quadratic
% lean: CausalSmith.Stat.WeakOverlap.lowerPair_bernoulli_kl_le_baseline_quadratic
% lean: CausalSmith.Stat.WeakOverlap.lowerPair_completionAt_observed_klDiv_le

\item Define the open bump box by
\[
E_h:=\{x\in\mathbb R^d:|x_i-x_{0i}|<h
\text{ for every }i=1,\ldots,d\}.
\]
Outside \(E_h\), at least one coordinate factor of the bump is zero.
Hence
\[
p_{0,h}(x)=p_{1,h}(x),
\qquad
\mathsf Q_{0,x}=\mathsf Q_{1,x},
\qquad x\notin E_h.
\]
For every \(r\ge0\), each truncated coordinate interval has length at
most \(2r\), so
\[
F_R(r)\le(2r)^d.
\]
For \(x\in E_h\), we have \(R(x)\le h\), and thus
\[
e_0(x)
=F_R(R(x))^{1/(\gamma-1)}
\le \bigl((2R(x))^d\bigr)^{1/(\gamma-1)}
\le \bigl((2h)^d\bigr)^{1/(\gamma-1)}.
\]
Moreover,
\[
\operatorname{vol}(E_h\cap\mathcal X)
\le\mathbb P\{R(X)\le h\}
=F_R(h)
\le(2h)^d.
\]

Both observed laws are obtained by integrating \(\mathsf Q_{j,x}\)
against the same uniform covariate law, and each conditional law
records its covariate \(x\) in the first coordinate. The entropy
disintegration identity therefore gives
\[
\operatorname{KL}(P_{0,h},P_{1,h})
=
\int_{\mathcal X}
\operatorname{KL}(\mathsf Q_{0,x},\mathsf Q_{1,x})\,dx.
\]
On \(E_h\), apply the conditional entropy and propensity bounds; on
its complement, the integrand is zero. Consequently,
\[
\begin{aligned}
\operatorname{KL}(P_{0,h},P_{1,h})
&\le
\frac{2}{\theta_0}(ah^\beta)^2
\bigl((2h)^d\bigr)^{1/(\gamma-1)}
\operatorname{vol}(E_h\cap\mathcal X)\\
&\le
\frac{2}{\theta_0}(ah^\beta)^2
\bigl((2h)^d\bigr)^{1/(\gamma-1)}(2h)^d\\
&=
\frac{2}{\theta_0}\,2^{D_\gamma}a^2h^{2\beta+D_\gamma}
=K_{\mathrm{KL}}a^2h^{\tau_{\mathrm{KL}}}.
\end{aligned}
\]
The last equality uses
\[
\frac{d}{\gamma-1}+d=D_\gamma.
\]
The finite entropy bound supplies both
\[
P_{0,h}\ll P_{1,h},
\qquad
\int
\left|\log\frac{dP_{0,h}}{dP_{1,h}}\right|\,dP_{0,h}<\infty.
\]

Finally, the oracle mesh satisfies
\[
h_{\gamma,n}^{\tau_{\mathrm{KL}}}=n^{-1},
\qquad
nh^{\tau_{\mathrm{KL}}}
=n(ch_{\gamma,n})^{\tau_{\mathrm{KL}}}
=c^{\tau_{\mathrm{KL}}}.
\]
The fixed calibration from the second step therefore gives
\[
n\operatorname{KL}(P_{0,h},P_{1,h})
\le K_{\mathrm{KL}}a^2c^{\tau_{\mathrm{KL}}}
\le\frac18.
\]
% lean: CausalSmith.Stat.WeakOverlap.lowerPair_observed_klDiv_le_raw
% lean: CausalSmith.Stat.WeakOverlap.lowerPair_raw_rate_identity
% lean: CausalSmith.Stat.WeakOverlap.lowerPair_observed_guards_and_kl_rate
% lean: CausalSmith.Stat.WeakOverlap.lowerPair_oracle_kl_balance

\item It remains to identify the chosen response representatives at
\(x_0\). The constructed memberships have the continuous Hölder
representatives \(\mu_{j,h}=Bp_{j,h}\). By
\cref{obj:lem:causal-identification}, both these representatives and the
chosen representatives associated with the same observed laws agree
almost everywhere with the treated conditional regression. Its
pointwise uniqueness conclusion gives
\[
\mu_{1,P_{j,h}}(x)=\mu_{j,h}(x)=Bp_{j,h}(x),
\qquad x\in\mathcal X,\quad j\in\{0,1\}.
\]
This includes boundary points of the cube.

Every coordinate factor of \(\psi\) equals one at zero, so
\[
\psi(0)=1,
\qquad
p_{1,h}(x_0)-p_{0,h}(x_0)=ah^\beta.
\]
Writing
\[
r_{\gamma,n}:=h_{\gamma,n}^{\beta},
\]
we conclude that
\[
\begin{aligned}
\left|\mu_{1,P_{1,h}}(x_0)-\mu_{1,P_{0,h}}(x_0)\right|
&=\left|Ba(ch_{\gamma,n})^\beta\right|\\
&=Bac^\beta r_{\gamma,n}
=c_0r_{\gamma,n}.
\end{aligned}
\]
The positivity of \(B,a,c\) justifies the second equality. Since \(n\ge1\)
was arbitrary, the fixed constants \(a,c,c_0\) satisfy all the stated
requirements.
% lean: CausalSmith.Stat.WeakOverlap.responseOf_eq_given_model_response
% lean: CausalSmith.Stat.WeakOverlap.pairSuccess_oracle_separation
% lean: CausalSmith.Stat.WeakOverlap.boundedLowerPair_spec
\end{enumerate}
\end{proof}
